\section{Related Work}\label{sec:related}

\subsection{Vision-Language-Action Models}
\textbf{Foundation VLA.} 
Vision-language-action foundation models typically adopt a powerful pre-trained vision-language model~\cite{bai2025qwen2, paligemma} as the semantic backbone, coupled with diffusion-based action head.
% By employing flow-matching processes, these models can generate continuous action chunks rather than single-step discrete tokens.
Recent VLA foundation models~\cite{pi_0, pi_0d5, gr00tN1, gr00t_n1_6, gemini_robotics, walloss, galaxeaG0, magma, gr3} have demonstrated enhanced multi-task execution capabilities and superior multi-embodiment adaptability, following pre-training on larger-scale and increasingly diverse datasets.
Distinguishing itself from the datasets utilized in preceding VLA foundation models, our model is pre-trained on an extensive corpus approximate 20,000 hours of multi-embodiment data. This massive-scale dataset, characterized by its high behavioral diversity, significantly bolsters the model's generalization capabilities across various robotic manipulation tasks.


\noindent
\textbf{Spatial VLA.} 
While traditional VLA models excel at semantic understanding, they often struggle with precise geometric reasoning and depth perception required for complex spatial manipulation.
To address this, several works~\cite{spatialvla, Spatial_forcing, internvla_M1, guo2025omnivla, magma, gemini_robotics_1d5, geovla} have integrated spatial representations into the VLA framework.
Several research~\cite{magma, internvla_M1, gemini_robotics} initiatives have focused on bolstering the spatial awareness of VLMs within embodied scenarios to enhance the spatial manipulation capabilities of VLAs in downstream tasks. 
Others explicitly or implicitly incorporate depth information during the VLA training phase.
Spatial Forcing~\cite{Spatial_forcing} employs a streamlined alignment strategy that compels the integration of VLA visual embeddings with spatial representations, thereby significantly improving the model's spatial comprehension.
% Different with the prior spatial VLA works, our method....\todo{}


\subsection{Evaluation on Robot Policy}
Current evaluation methodologies for robot policies are primarily bifurcated into two categories: simulation-based~\cite{libero, calvin, simplerenv, robocasa, robotwin2} and real-world embodiment-based~\cite{robochallenge, roboarena}.
Simulation-based benchmark provide a rapid and convenient means to evaluate the capabilities of policies, enabling large-scale parallel testing across vast and diverse interaction scenarios at very low cost. 
Although simulation environments typically employ idealized physical models, their results often do not fully represent the complexity of the real physical world.
The another real-world evaluations' efficiency is often bottlenecked by the requirement for extensive hardware parallelism. 
Consequently, the majority of prior VLA studies have been confined to comparing a limited number of methods across only a few tasks. 
% Recently, RoboChallenge~\cite{robochallenge} introduced a framework for the fair comparison of multiple policies across 30 tasks via remote robotic arm deployment.
To more comprehensively evaluate the real-world performance of policies, this work conducts an assessment across three distinct robotic platforms, with 100 tasks executed on each platform. We further provide a thorough analysis of how mainstream VLA models adapt to the diversity encountered in real-world scenarios.


\subsection{Efficient VLA Training}
The rapid iteration of VLA models has catalyzed the development of specialized training infrastructure. Several well-designed open-source codebases have recently emerged in the community, each catering to different research priorities. 
For instance, the OpenPI~\cite{pi_0} repository provides a versatile framework supporting both JAX and PyTorch for training the $\pi$ series models. 
StarVLA~\cite{starvla} introduces a modular and user-friendly codebase specifically optimized for the co-training of VLAs and VLMs, facilitating the transfer of semantic knowledge to robotic control. 
Additionally, Dexbotic~\cite{dexbotic} is designed as a unified and efficient solution to streamline the development lifecycle of VLAs, focusing on standardizing the pipeline from data ingestion to model deployment.
Despite these advancements, training large-scale VLA models on multi-node clusters remains a significant challenge due to data I/O bottlenecks and communication overheads.
To bridge this gap, we present \method, a high-performance open-source codebase engineered for large-scale VLA training. Unlike existing frameworks, our codebase implements systemic optimizations in data loading, distributed training strategies, and operator-level acceleration. 
These enhancements lead to a comprehensive improvement in training throughput and scalability, providing a more efficient foundation for the community to explore the scaling limits of robotic foundation models.




% \lipsum[1-7]
